% !TeX program = lualatex
% !TeX encoding = utf8
% !TeX root = ../Thermod.tex

%\setcounter{chapter}{13}

%=========================================================
\Opensolutionfile{answer}[\currfilebase/\currfilebase-Answers]
\Writetofile{answer}{\protect\section*{\nameref*{\currfilebase}}}
\chapter{Элементи кінетики і процеси переносу}\label{\currfilebase}
\graphicspath{{\currfilebase/pictures}}
\makeatletter
\def\input@path{{\currfilebase/}}
\makeatother
%=========================================================

%% --------------------------------------------------------
\section{Довжина вільного пробігу}
%% --------------------------------------------------------

Ключову роль в молекулярній теорії газів, зокрема, в питаннях кінетики відіграє
поняття середньої довжини вільного пробігу~$l$. Якщо вважати молекули газу
точковими об'єктами, то зіткнення між ними неможливі, а отже, неможливе
встановлення теплової рівноваги. Однак, взаємодія реальних молекули є силами
тяжіння $\sim 1/r^6$ на великих відстанях і відштовхування $\sim 1/r^{12}$ на
малих (такими, що можна порівняти з розмірами молекул). Область найбільшої
взаємодії має порядок розмірів молекул і залежить від відносної швидкості руху
молекул, а отже, і від температури. У першому наближенні цією властивістю можна
знехтувати. Будемо розглядати молекули як тверді кулі діаметром~$d$.

Під час руху молекули між двома зіткненнями вона вимітає об’єм, що дорівнює
об'єму циліндра $\pi d^2 l$, недоступний іншим молекулам
(рис.~\ref{pic:Pic15.1}). Дійсно, дві молекули можуть зіткнутися, якщо
відстань між їх центрами менше ніж~$d$. Середній об’єм циліндра, що вимітається,
дорівнює $\pi d^2 l$, де $l$ і є середня довжина вільного пробігу. На цей об’єм
припадає одна молекула, а це означає, що концентрація молекул газу може бути
визначена як

\begin{equation}
	n = \frac{1}{\pi d^2 l}.
	\label{eq:150.1}
\end{equation}

Звідси

\begin{equation*}
	l = \frac{1}{\pi d^2 n}.
\end{equation*}
%
де $\sigma$~--- ефективний газо-кінетичний перетин зіткнення, а $d$~---
газо-кінетичний діаметр молекули. Насправді, користуючись
формулою~\eqref{eq:150.1} якраз і можна знайти діаметр молекул, точніше
ефективний діаметр, скориставшись експериментальними значеннями $l$ і $n$.
Слово «ефективний» підкреслює метод визначення цього діаметру. Зрозуміло, що
можна скористатися його значенням, отриманим із інших методів.

%---------------------------------------------------------
\begin{figure}[h!]\centering
	\includegraphics[width=0.75\linewidth]{Pic1}
	\caption{До визначення середньої довжини вільного пробігу.}
	\label{pic:Pic15.1}
\end{figure}
%---------------------------------------------------------

Так, для \ce{He} і кількості атомів при $0^\circ$C і тиску 760 мм рт.~ст. в
$1$~см$^3$ концентрація становить
$n = 2.69\cdot 10^{19}\,\text{см}^{-3}$. Звідки отримаємо:

\begin{equation*}
	l = 2.5 \cdot 10^{-5}\,\text{см} = 2500\ \text{\AA}.
\end{equation*}

Корисно оцінити і частоту зіткнень як:

\begin{equation}
	\nu = \frac{1}{2} n \sigma \langle v \rangle,
	\label{eq:150.2}
\end{equation}

де множник $1/2$ враховує, що в зіткненнях бере участь в основному дві молекули.
Тоді для \ce{He} отримаємо, що $\nu \approx 10^{21}$ зіткнень/(сек·см$^3$).

%% --------------------------------------------------------
\section{Поняття ефективного перетину}
%% --------------------------------------------------------

Слід зазначити, що формула~\eqref{eq:150.1} не є точною. Якщо врахувати, що
молекули рухаються і прийняти, що їх швидкості підкоряються розподілу Максвелла,
то довжина вільного пробігу виявляється в $\sqrt{2}$ разів меншою. Однак такого
роду уточнення, хоча і робилися такими геніями, як Максвелл, є малоцінними,
оскільки саме поняття «діаметр» стосовно до реальної молекули є досить
примарним.

Навпаки, поняття ефективного перетину~$\sigma$ набуло широкого поширення в
фізиці. Величина ця може бути визначена експериментально, наприклад, у дослідах
по зіткненню пучків молекул, атомів, елементарних частинок із мішенню, як
відношення середнього числа частинок~$\Delta N$, вибитих із пучків в одиницю
часу при зіткненнях, наслідком яких є необхідний результат, до інтенсивності
пучка:

\begin{equation*}
	\sigma = \frac{\Delta N}{I}.
\end{equation*}

Тому говорять про ефективні перетини іонізації, поглинання, розсіювання,
реакції та інше --- залежно від того, який процес розглядають. Відповідно, можна
говорити і про довжини вільного пробігу для кожного із цих процесів.

%% --------------------------------------------------------
\section{Імовірність зіткнень}
%% --------------------------------------------------------

Імовірність того, що зіткнення станеться на довжині $dL$ дорівнює
$\beta\, dL$, тобто вона пропорційна $dL$ і не залежить від уже пройденого
шляху. Позначимо $f(L)$ — ймовірність того, що після зіткнення молекула пройде
шлях~$L$. Імовірність того, що на наступному відрізку $dL$ зіткнення не
відбудеться, дорівнює $(1 - \beta\, dL)$.

Тоді ймовірність $f(L + dL)$ того, що молекула пройде без зіткнення дистанцію
$L + dL$, є добутком двох незалежних подій:
ймовірності того, що вже пройдено шлях $L$, і ймовірності того, що без
зіткнень буде пройдено ще $dL$:

\begin{equation}
	f(L + dL) = f(L)\,(1 - \beta\, dL).
	\label{eq:152.1}
\end{equation}

Звідки:

\begin{equation}
	\frac{f(L + dL) - f(L)}{dL} = -\beta\, f(L),
	\label{eq:152.2}
\end{equation}

або в границі $dL \to 0$:

\begin{equation}
	\frac{df}{dL} = -\beta f.
	\label{eq:152.3}
\end{equation}

Оскільки $f(0) = 1$, то $A = 1$ у розв’язку рівняння. Таким чином:

\begin{equation*}
	f(L) = e^{-\beta L}.
\end{equation*}

Щоб знайти $\beta$, визначимо середню відстань, яку проходить молекула між
зіткненнями. За означенням вона дорівнює довжині вільного пробігу:

\begin{equation}
	l = \int_0^\infty L\, f(L)\, dL.
	\label{eq:152.5}
\end{equation}

Тоді остаточно ймовірність того, що після зіткнення молекула пройде шлях $L$:

\begin{equation}
	f(L) = \frac{1}{l}\, e^{-L/l}.
	\label{eq:152.6}
\end{equation}

%---------------------------------------------------------
\begin{figure}[h!]\centering
	\includegraphics[width=0.75\linewidth]{Pic2}
	\caption{Стаціонарний тепло- (а) та масо- (б) перенос між двома системами}
	\label{pic:Pic152.1}
\end{figure}
%---------------------------------------------------------

% Друга ілюстрація:
\begin{figure}[h!]\centering
	\includegraphics[width=0.5\linewidth]{Pic3}
	\caption{Два режими руху молекул і два режими переносу.}
	\label{pic:Pic152.2}
\end{figure}
%---------------------------------------------------------

\section{Явища і режими переносу фізичних величин}

Розглянемо систему, яка не знаходиться в стані теплової рівноваги. Практичний
інтерес мають системи, що знаходяться в стаціонарних нерівноважних станах.
При цьому від однієї частини системи до другої переноситься якась фізична
величина — є нескомпенсований макроскопічний потік.

Найбільш очевидними кандидатами на такі фізичні величини є перенос енергії у
вигляді тепла або перенос кількості частинок (рис.~\ref{pic:Pic152.1}). При
цьому рушійною силою процесів переносу є перепад температур на одиницю довжини
(градієнт температури) для процесів теплопровідності або перепад концентрацій
для процесів переносу частинок відповідно.

Як теоретичні, так і експериментальні результати вказують на те, що між потоком
відповідної фізичної величини ($J_E$, $J_n$) і рушійною силою існує пряма
пропорційність. Значення коефіцієнта пропорційності залежить від режиму
переносу (рис.~\ref{pic:Pic152.2}).

Справа в тому, що фізичними носіями, які здійснюють процес переносу, є частинки
речовини (молекули, іони, електрони). Зазвичай довжина вільного пробігу
$l \ll L$ (де $L$ --- характерний розмір системи), і молекула зазнає величезної
кількості зіткнень, перш ніж подолає шлях від одного кінця системи до іншого.
В окремих ділянках системи на шляху руху молекули встановлюється локальна
термодинамічна рівновага, що характеризується локальною температурою,
концентрацією, тиском тощо.

Натомість за тиску газу $\sim 10^{-6}$ атм. або $\sim 10^{-3}$ мм рт. ст.,
довжина вільного пробігу стає $l \sim 10$~см, тобто порівнянною з розмірами
лабораторних установок, за допомогою яких ми досліджуємо процеси переносу.
У цьому випадку молекули здійснюють перенесення фізичних величин
безпосередньо, без проміжних зіткнень. Обрахунки стають принципово іншими,
ніж у випадку $l \ll L$.

%% --------------------------------------------------------
\section{Умови \texorpdfstring{$l \ll L$}{l << L}. Локальна термодинамічна рівновага}
%% --------------------------------------------------------

Розглянемо процеси переносу в умовах, з якими ми зазвичай маємо справу, тобто
коли $l \ll L$. Послідовність розрахунків однакова для всіх процесів переносу
незалежно від того, чи то процеси дифузії, теплопровідності, в’язкості або
електропровідності. Фізичними носіями для перших трьох процесів є молекули, якщо
мова йде про гази.

Рух молекул газу ізотропний, тому ми можемо скористатися результатами
молекулярно-кінетичної теорії, згідно з якою кількість частинок із швидкістю
$v$, які рухаються під кутом~$\theta$ до нормалі площини $zy$
(рис.~\ref{pic:Pic154.1}), що переносяться в одиницю часу через одиницю площі,
становить:

\begin{equation}
	\frac{1}{4}\, n_0\, \varphi(v)\, v\, dv\, \sin\theta\, d\theta.
	\label{eq:154.1}
\end{equation}

Продовжимо для конкретності розгляд процесу переносу енергії цими частинками.
Згідно з (152.6) кожна з цих молекул зазнає зіткнень, які спричиняють зміни
енергії частинок, на відстані $dx$ від площини $zy$ з координатою 0 з
ймовірністю $\exp\{-x/(l\cos\theta)\}\, dx/\cos\theta$, де $l$ — довжина
вільного пробігу, а $L = dx/\cos\theta$ — шлях, пройдений молекулою без
зіткнень під кутом~$\theta$ до нормалі площини $zy$.

Отже, на відстані $l \cos\theta$ від площини $zy$ вони в основному зазнають
зіткнення. Тому вважатимемо, що в кожний момент часу в площинах $zy$ з різними
$x$ координатами встигає встановитися термодинамічна рівновага. Будемо говорити
про неї як про \emph{локальну термодинамічну рівновагу}.

Енергія цих молекул у середньому дорівнює $\varepsilon(x)$ в площині $x$ і
$\varepsilon(x+dx)$ в площині $x+dx$. Тоді потік енергії від площадки з
координатою $x$ до площадки з координатою $x+dx$, що здійснюється частинками зі
швидкістю $v$, які рухаються під кутом $\theta$ до осі $x$, становить:

\begin{equation}
	dJ_E = \frac{1}{4} n_0\, \varphi(v)\, v\, dv\, \sin\theta\, d\theta\,
	\bigl[ \varepsilon(x) - \varepsilon(x+dx) \bigr]\,
	\exp\{-x/(l\cos\theta)\}\, .
	\label{eq:154.2}
\end{equation}

Розкладаючи в ряд
\[
	\varepsilon(x+dx)=\varepsilon(x)+\frac{d\varepsilon}{dx}\,dx,
\]
вираз~\eqref{eq:154.2} може бути переписаний як:

\begin{equation}
	dJ_E = -\frac{1}{4} n_0\, \varphi(v)\, v\, dv\, \sin\theta\, d\theta\,
	\frac{d\varepsilon}{dx}\, dx.
	\label{eq:154.3}
\end{equation}

Повний потік можна отримати інтегруванням~\eqref{eq:154.3} по відстані, яку
проходить молекула під кутом~$\theta$ до нормалі, $dL = dx/\cos\theta$, від 0 до
$l\cos\theta$, по $d\theta$ від $0$ до $\pi$, і по $dv$ від $0$ до~$\infty$:

\begin{equation}
	J_E = -\frac{1}{4}\, \frac{d\varepsilon}{dx}
	\int_0^\infty n_0 \varphi(v)\, v\, dv
	\int_0^\pi \sin\theta\, \cos\theta\, d\theta.
	\label{eq:154.4}
\end{equation}

Інтегрування від 0 до $\pi$ враховує молекули, які йдуть з площини $zy$ з
координатою $x$, та молекули, що приходять до неї знизу. Враховуючи, що
$T = T(x)$, отримаємо:

\begin{equation}
	J_E = -\frac{1}{3}\, n_0\, \langle v \rangle\, l\,
	\frac{d\varepsilon}{dx}.
	\label{eq:154.5}
\end{equation}

Оскільки $n_0 \varepsilon = U/V$, то остаточно:

\begin{equation}
	J_E = -\frac{1}{3} C_V\, \langle v \rangle\, l\,
	\frac{dT}{dx},
	\label{eq:154.6}
\end{equation}
%
де $\widehat{C}_V$ --- теплоємність одиниці об’єму при сталому об’ємі $V$.

Отриманий закон називають \emph{першим законом Фур'є}, який пов’язує потік
енергії з градієнтом температури. Коефіцієнт пропорційності називається
\emph{коефіцієнтом теплопровідності}~$\lambda$:

\begin{equation}
	J_E = -\lambda\, \frac{dT}{dx},
	\qquad
	\lambda = \frac{1}{3}\, \widehat{C}_V\, \langle v \rangle\, l.
	\label{eq:154.7}
\end{equation}

В цілому, чим важчий газ, тим гірше його теплопровідність при заданій
температурі та тиску, оскільки менша середня швидкість молекул, а значення
теплоємності одиниці об’єму газу залежить лише від кількості ступенів вільності
молекул (у моделі ідеального газу).

%---------------------------------------------------------
\begin{figure}[h!]\centering
	\includegraphics[width=0.5\linewidth]{Pic4}
	\caption{До виводу першого закону Фур'є}
	\label{pic:Pic154.1}
\end{figure}
%---------------------------------------------------------

%% --------------------------------------------------------
\section{Перший закон Фіка}
%% --------------------------------------------------------

Схожим чином можна знайти коефіцієнт дифузії. Рушійною силою цього процесу буде
градієнт концентрації $n(x)$ частинок. Зіткнення, які перешкоджають вільному руху
частинок, також відбуваються із певною ймовірністю і в основному стаються після
проходження молекулами шляху довжиною $l$.

Тоді потік частинок зі швидкістю $v$, які рухаються під кутом $\theta$ до осі $x$
від площадки з координатою $x$ до площадки з координатою $x+dx$, буде
(по аналогії з ходом міркувань (154.2–154.6)):

\begin{equation}
	dJ_n = -\frac{1}{4}\, n_0\, \varphi(v)\, v\, dv\, \sin\theta\, d\theta\,
	\frac{dn}{dx}\, dx.
	\label{eq:155.1}
\end{equation}

Інтегруючи~\eqref{eq:155.1} по всіх можливих напрямках і швидкостях:

\begin{equation}
	J_n = -\frac{1}{3}\, \langle v \rangle\, l\, \frac{dn}{dx}.
	\label{eq:155.2}
\end{equation}

Звідси коефіцієнт дифузії:

\begin{equation}
	D = \frac{1}{3}\, \langle v \rangle\, l.
	\label{eq:155.3}
\end{equation}

Закон~\eqref{eq:155.2} отримав назву \emph{першого закону Фіка}.

Отриманий вираз дає значення коефіцієнта \emph{самодифузії} молекул того ж самого
сорту. Очевидно, що коефіцієнт дифузії молекул сорту A в газі B буде відрізнятися
через інші значення довжини вільного пробігу і характер взаємодії молекул різних
газів. Урахування цих обставин є предметом окремої дискусії.

Для того, щоб визначити коефіцієнт самодифузії, слід знайти спосіб “прослідкувати”
певну молекулу. На практиці це робиться або реєстрацією дифузійного потоку
ізотопів молекул, або молекул близьких за масою і ефективним перерізом зіткнень.
Якщо в початковий момент часу ці молекули розділені у просторі, то з часом
утворюються взаємно спрямовані потоки молекул різного сорту:
\[
	J_{n_1} = - J_{n_2}.
\]
Останнє випливає з локальної термодинамічної рівноваги, яка характеризується
умовою $n_1(x)+n_2(x)=n_0$. Продиференціювавши, отримаємо умову рівності потоків.

%% --------------------------------------------------------
\section{Формула Ейнштейна}
%% --------------------------------------------------------

Уявімо собі <<мічену>> частинку, яка рухається хаотично, зазнаючи постійних
зіткнень із іншими молекулами. Її можна розглядати як браунівську частинку.
За відсутності зовнішніх полів її положення змінюється випадково. Проте за
наявності зовнішнього поля (наприклад, гравітаційного або електричного, якщо
частинка заряджена) з’являється нескомпенсована швидкість $\vect{u}$.

Експеримент свідчить, , що при відносно невеликих значеннях швидкості $\vect{u}$, між нею і силою $\vect{F}$, що на неї діє, існує проста залежність:

\begin{equation}
	u = B F,
	\label{eq:155.4}
\end{equation}
%
де $B$ --- рухливість частинки.

Під дією цієї сили встановлюється певний концентраційний розподіл $n(x)$, який
описується больцманівським розподілом:

\begin{equation}
	n(x) = n_0\, \exp\!\left(-\frac{F x}{k T}\right).
	\label{eq:155.5}
\end{equation}

Але цей розподіл спричиняє появу дифузійного потоку, який намагається вирівняти
концентрацію. В рівновазі він компенсується потоком $n\vect{u}$, зумовленим дією сили:

\begin{equation}
	- D\, \frac{dn}{dx} + n u = 0.
	\label{eq:155.6}
\end{equation}

Розв’язуючи рівняння та порівнюючи з~\eqref{eq:155.5}, отримаємо зв’язок між
коефіцієнтом дифузії та рухливістю:

\begin{equation}
	D = B\, k T.
	\label{eq:155.7}
\end{equation}

З іншого боку рух частинки маси $m$ під дією сили $F$, що зазнає випадкових
поштовхів від інших молекул $\xi(t)$, описується рівнянням:

\begin{equation}
	m \frac{d^2 x}{dt^2} = F + \xi(t) - \alpha \frac{dx}{dt},
	\label{eq:155.8}
\end{equation}

де $\alpha$ — коефіцієнт опору.

Домноживши рівняння на $x$ та усереднивши:

\begin{equation}
	m \frac{d^2}{dt^2} \langle x^2 \rangle
	= 2 k T - 2 \alpha \frac{d}{dt} \langle x^2 \rangle .
	\label{eq:155.9}
\end{equation}

Зручніше розв’язувати рівняння для похідної:

\begin{equation}
	\frac{d}{dt}\langle x^2 \rangle
	= \frac{2 k T}{\alpha}
	\left(1 - e^{-2\alpha t / m}\right).
	\label{eq:155.10}
\end{equation}

У границі $t \to \infty$ середній квадрат зміщення:

\begin{equation}
	\langle x^2 \rangle = 2 D t,
	\label{eq:155.11}
\end{equation}
%
з урахуванням формули Ейнштейна $D = kT/\alpha$.

Таким чином \emph{середній квадрат зміщення частинки в певному напрямку зумовлений дифузією пропорційний часу спостереження за нею}. Вираз, що отримав назву \emph{формули Ейнштейна}, дає можливість визначити по цій залежності коефіцієнт дифузії і в тому чи іншому вигляді лежить в основі всіх експериментальних методів.

\section{Коефіцієнт в’язкості}

Для того, щоб отримати вираз для коефіцієнта в’язкості, слід визначити, яка
фізична величина переноситься в такому процесі. Розглянемо ламінарну течію газу
(або рідини) між двома площинами: рухомою aa' і нерухомою bb' (Рис.~\ref{pic:Pic157.1}).
Крім молекулярної швидкості, яка визначається температурою газу, молекули
будуть мати і макроскопічну компоненту швидкості, яка направлена уздовж у і
змінюється за абсолютною величиною уздовж х. Причина появи цією
нескомпенсованої компоненти є сили в’язкого тертя, що діють між різними шарами
газу або рідини. Це можна трактувати так, що компоненти швидкості $v_y$
молекули змінюється при переході з площини zy з координатою х в площину
з координатою $x+dx$.

Таким чином в даному випадку молекули переносять компоненту швидкості $v_y$
і виникає потік y-вої компоненти імпульсу. Отже молекули рухаючись вгору
гальмують рух верхніх шарів, тому що мають менший імпульс і навпаки.
Міркуючи у такий спосіб, як і раніше знайдемо, що потік компоненти імпульсу
$p_y$, що переноситься частинками зі швидкістю $v$, які рухаються під кутом
$\theta$ до осі х від площадки з координатою x до площадки з координатою
$x+dx$:

\begin{equation}
	dJ_{p_y} = -\frac{1}{4} n_0 \varphi(v) v \sin\theta \, dv \, d\theta \,
	\frac{dp_y}{dx} \, dx .
	\label{eq:157.1}
\end{equation}

Тоді повний потік уздовж х знаходиться інтегруванням:

\begin{equation}
	J_{p_y} = -\frac{1}{3} \langle v \rangle l \frac{dp_y}{dx}.
	\label{eq:157.2}
\end{equation}

Коефіцієнт:

\begin{equation}
	\eta = \frac{1}{3} \rho \langle v \rangle l .
	\label{eq:157.3}
\end{equation}
%
називають коефіцієнтом в’язкості. Відзначимо, що $\eta$ в газах не залежить
від тиску тому що $l \sim \rho^{-1}$, що було перевірено експериментально аж до
тиску $1/60$ атм. Цей висновок був з одним перших переконливих доказів
в справедливості молекулярно-кінетичної теорії. Зауважимо, що з точки зору
механіки потік $J_{p_y}$ є за визначенням відношенням
$d p_y / d t\, d\Sigma_{yz} = \sigma_{yx}$, де $d\Sigma_{yz}=A$ --- ділянка одиничної
площі в площині zy, а $dp_y/dt$ є сила прикладена до цієї площадки (Рис.~\ref{pic:Pic157.1}б).
Тоді $\sigma_{yx}$ є напруження зсуву прикладене до ділянки $d\Sigma_{yz}$
і \eqref{eq:157.2} переписується так:

\begin{equation}
	\sigma_{yx} = \eta \frac{d v_y}{dx}.
	\label{eq:157.4}
\end{equation}
%
в повній згоді з гідродинамічним трактуванням поняття в'язкості, згідно із яким
напруження зсуву пропорційне градієнту швидкості.

%---------------------------------------------------------
\begin{figure}[h!]\centering
	\includegraphics[width=0.95\linewidth]{Pic5}
	\caption{До виведення формули для коефіцієнту в’язкості в термодинаміці (а)
		і поняття в’язкості з точки зору механіки (б)}
	\label{pic:Pic157.1}
\end{figure}
%---------------------------------------------------------

%% --------------------------------------------------------
\section{Кінетичні коефіцієнти}
%% --------------------------------------------------------

Зручно подати отримані результати у вигляді таблиці, узагальнивши їх і на
явище електропровідності, яке буде розглянуто згодом.

\begin{table}[h!]
	\centering\small
	\begin{tblr}{
			colspec={Q[c,m, 2.5cm] Q Q[c,m] Q[c,m] Q[c,m] Q[c,m] },
			hlines,
			vlines,
			row{1} = {font=\bfseries, c, m, bg=cyan!50}
		}
		Процес                        & Носій & Потік & Градієнт &
		Коефіцієнт                    & Закон                                     \\

		Дифузія                       &
		Частинки                      &
		$J_n$                         &
		$\nabla n$                    &
		$D$                           &
		$J_n = - D \, \nabla n$                                                   \\

		Тепло\-провідність              &
		{Молекули,\\ електрони}       &
		$J_E$                         &
		$\nabla T$                    &
		$\lambda$                     &
		$J_E = - \lambda \, \nabla T$                                             \\

		В’язкість                     &
		Молекули                      &
		{Потік\\ імпульсу:\\ $\sigma_{yx}$} &
		$\dfrac{dv_y}{dx}$            &
		$\eta$                        &
		$\sigma_{yx} = \eta \dfrac{dv_y}{dx}$                                     \\

		Електро\-провідність            &
		{Електрони,\\ іони}           &
		$\mathbf{j}$                  &
		{Електричне\\ поле\\ $\mathbf{E} = - \nabla\phi$}  &
		$\sigma$                      &
		$\mathbf{j} = \sigma \mathbf{E}$                                          \\
	\end{tblr}
\end{table}

Співставлення виразів отриманих для різних газокінетичних коефіцієнтів
дозволяє знайти простий зв'язок між ними. Так коефіцієнт теплопровідності і
дифузії пов’язані простим співвідношенням:

\begin{equation}
	\lambda = C_V D ,
	\label{eq:157.1}
\end{equation}

Далі зауважимо, що $\hat{C}_V = c_v \rho$, де $c_v$ — питома теплоємність.
Тоді знаходимо, що:

\begin{equation}
	\lambda = c_v \rho D ,
	\label{eq:157.2}
\end{equation}

Це співвідношення неодноразово перевірялось експериментально і було знайдено,
що швидше виконується співвідношення
$\lambda = f \cdot \hat{C}_V \cdot \eta$, де $f$ змінюється від $1.3$ до $2.3$.
Наявність числового множника пов'язане зі спрощеною моделлю, якою ми
користувалися, не враховуючи зокрема того, що перенесення обертальної і
поступальної енергії відбувається по різному, а також того, що усереднення
проводилося досить грубо. Теоретичний розрахунок, проведений Джинсом,
дозволяє знайти цей коефіцієнт $f$.

Незважаючи на зазначені неточності, вирази для коефіцієнтів переносу дають
можливість визначити функціональні залежності від температури, тиску та
молекулярної ваги, принаймні для газів, для яких наближення
молекулярно-кінетичної теорії цілком виправдані. Для оцінки величин, що
входять в формули для коефіцієнтів переносу будемо вважати, що гази
задовольняють рівнянню Клайперона–Менделєєва, мають розподіл Максвела і
довжина вільного пробігу зворотно пропорційна концентрації молекул газу:

\begin{equation}
	l \sim \frac{1}{n \sigma},
	\label{eq:157.3}
\end{equation}

Використовуючи отримані співвідношення, знайдемо для коефіцієнтів переносу:

\begin{equation}
	D \sim \frac{1}{3} \langle v \rangle l,
	\qquad
	\eta \sim \frac{1}{3} \rho \langle v \rangle l,
	\qquad
	\lambda \sim \frac{1}{3} \hat{C}_V \langle v \rangle l .
	\label{eq:157.4}
\end{equation}

Відзначимо ще одну обставину. Виміри коефіцієнта в'язкості і теплопровідності
дозволяють експериментально знайти значення довжини вільного пробігу, а відтак
і діаметр молекул. При цьому значення, отримані з коефіцієнтів в'язкості і
теплопровідності, відрізняються незначним чином. Так, для H$_2$:
$d(\lambda) = 2{,}40\,\text{\AA}$, а $d(\eta) = 2{,}47\,\text{\AA}$.
З урахуванням простоти засобів, якими були отримані дані результати, можна
сказати, що це блискучий результат.

%% --------------------------------------------------------
\section{Закон Відемана–Франца}
%% --------------------------------------------------------

Причина, з якої існують настільки прості співвідношення між кінетичними
коефіцієнтами очевидна --- фізичними носіями у всіх процесах є частинки газу,
які зазнають випадкові поштовхи з боку інших молекул і знаходяться в
хаотичному русі. Подібна ситуація спостерігається і в металах, для яких
експериментально було встановлено закон Відемана–Франца: відношення
коефіцієнта теплопровідності до коефіцієнта електропровідності має одне і
те ж значення для всіх металів при одній і тій же температурі.

Зрозуміло, що фізичними носіями в металах можуть бути тільки електрони.
Ця ситуація наштовхує на думку, що якщо електрони деякою мірою можна
розглядати як ідеальний газ із внутрішньою енергією
$U = \tfrac{3}{2}RT$, то сам процес електропровідності може бути описаний
в рамках молекулярно-кінетичних уявлень. Якщо це дійсно так, коефіцієнт
теплопровідності металу має дорівнювати:

\begin{equation}
	\lambda = C_V D ,
	\label{eq:158.1}
\end{equation}

%---------------------------------------------------------
\begin{figure}[h!]\centering
	\includegraphics[width=0.5\linewidth]{Pic6}
	\caption{До виведення значення коефіцієнту електропровідності.
		Електрони під дією електричного поля починають прискорюватися.
		Проте зіткнення із дефектами і збурення іонів кристалічної решітки
		гальмують рух. В результаті електрони дрейфують із сталою швидкістю
		проти прикладеного поля.}
	\label{pic:Pic6}
\end{figure}
%---------------------------------------------------------

Для обчислення коефіцієнта електропровідності, зауважимо, що на електрон
в електричному полі $E$ діє сила $eE$. Під дією цієї сили електрон
``дрейфує'' проти напряму поля з деякою кінцевою швидкістю $u$ (Рис.~\ref{pic:Pic6}).
Дійсно, компонента швидкості паралельна напрямку поля, яку набуває електрон
під дією сили $eE$ в середньому дорівнює:

\begin{equation}
	u = \frac{eE\tau}{m},
	\label{eq:158.2}
\end{equation}

де $\tau$ — час вільного пробігу електрона. Він може бути оцінений як
$\tau = l / \langle v \rangle$, де $l$ — довжина вільного пробігу електрона.
Звідси швидкість дрейфу $u$ дорівнює:

\begin{equation}
	u = \frac{eEl}{m\langle v \rangle}.
	\label{eq:158.3}
\end{equation}

За визначенням густина струму $J$ дорівнює кількості електрики, яка пройшла
через одиничну площадку в одиницю часу:

\begin{equation}
	J = ne\,u ,
	\label{eq:158.4}
\end{equation}

де $n$ — концентрація вільних електронів. Підставивши \eqref{eq:158.3}
у \eqref{eq:158.4}, одержимо закон Ома в диференційній формі
$J = \sigma E$ і вираз для питомої провідності:

\begin{equation}
	\sigma = \frac{n e^{2} \tau}{m}
	= \frac{n e^{2} l}{m \langle v \rangle}.
	\label{eq:158.5}
\end{equation}

Якщо взяти відношення $\lambda / \sigma$, знайдемо:

\begin{equation}
	\frac{\lambda}{\sigma} = \frac{C_V D}{\sigma}.
	\label{eq:158.6}
\end{equation}
%
де з точністю, припустимою в даній моделі, можна вважати, що
$m\langle v \rangle^{2}/2 \approx m\langle v^{2} \rangle/2 \approx \tfrac{3}{2}k_B T$.

Співвідношення \eqref{eq:158.6} і є закон Відемана–Франца. Коефіцієнт
пропорційності, що стоїть перед температурою, отримав назву
**число Лоренца**. Якщо $\lambda$ виміряти у
Вт/(см·град), а $\sigma$ в $\Omega^{-1}/\text{см}^3$, то значення
$\lambda / (\sigma T)$ при кімнатній температурі близьке до
$3(k/e)^2 = 2.44 \cdot 10^{8}$. Для різних металів воно змінюється
в межах приблизно $(2 \div 3)\cdot10^{8}$.

При зниженні температури спостерігається суттєва розбіжність між
теорією і експериментом, що зумовлено тим, що електронний газ при
низьких температурах вже не можна коректно описувати як класичний
ідеальний газ --- квантові ефекти стають важливими.

%% --------------------------------------------------------
\section{Другий закон Фур’є}
%% --------------------------------------------------------

%---------------------------------------------------------
\begin{figure}[h!]\centering
	\includegraphics[width=0.5\linewidth]{Pic7}
	\caption{До виводу рівняння теплопровідності (другого закону Фурьє).}
	\label{pic:Pic7}
\end{figure}
%---------------------------------------------------------

Тепер розглянемо деякий нестаціонарні задачі переносу. Принципи розв’язку задач такого роду вивчимо на прикладі розв’язку рівняння теплопровідності. Основи математичної теорії теплопровідності були закладені Жан Батістом Жозефом Фур'є. Хоча він і виходив із уявлень хибної теорії теплороду, подання тепла як певної рідини дозволяє звести процеси його поширення тепла до процесів току рідини, що не стискається. Різниця між потоками тепла через деякий одиничний перетин~$\Sigma$ з координатами $x$ і $x+dx$, дорівнює кількості тепла, яке накопичуються або виходить з об’єму між двома площинами (Рис.~\ref{pic:Pic7}):

\begin{equation}
	J(x) - J(x+dx) = \rho c_V \frac{\partial T}{\partial t}\,dx
	\label{eq:159.1}
\end{equation}

Поділивши \eqref{eq:159.1} на $dx$ і спрямувавши до нуля отримаємо:

\begin{equation}
	-\frac{dJ}{dx} = \rho c_V \frac{\partial T}{\partial t}
	\label{eq:159.2}
\end{equation}

Використовуючи перший закон Фур'є знайдемо далі:

\begin{equation}
	J = -\lambda \frac{\partial T}{\partial x}
	\qquad\Rightarrow\qquad
	\lambda\,\frac{\partial^2 T}{\partial x^2}
	= \rho c_V \frac{\partial T}{\partial t}
	\label{eq:159.3}
\end{equation}

Отримане рівняння~\eqref{eq:159.3} називають другим законом Фур’є або рівнянням теплопровідності. Якщо в середовищі діють внутрішні джерела тепла, що виділяють $q$ одиниць тепла в одиницю часу в одиниці об'єму, то \eqref{eq:159.3} переписується як:

\begin{equation}
	\lambda\,\frac{\partial^2 T}{\partial x^2} + q
	= \rho c_V \frac{\partial T}{\partial t}
	\label{eq:159.4}
\end{equation}

Зрештою для однорідного середовища, в якій коефіцієнт теплопровідності не залежить від температури знаходимо:

\begin{equation}
	\frac{\partial T}{\partial t}
	= \chi\,\frac{\partial^2 T}{\partial x^2},
	\qquad
	\chi=\frac{\lambda}{\rho c_V}
	\label{eq:159.5}
\end{equation}

де $\chi = \lambda/(\rho c_V)$ --- коефіцієнт температуропровідності.
Вирази \eqref{eq:159.3} --- \eqref{eq:159.5} були отримані для одновимірного потоку тепла.
У загальному випадку, коли властивості і температура середовища залежать від всіх трьох координат середовища $x, y, z$, рівняння \eqref{eq:159.4} записується так:

\begin{equation}
	\frac{\partial T}{\partial t} = \chi \nabla^2 T
	\label{eq:159.6}
\end{equation}

де вираз для градієнта визначається симетрією задачі (Рис.~\ref{pic:Pic8}). Зокрема, в декартовій системі координат:

\begin{equation}
	\nabla^2 T =
	\frac{\partial^2 T}{\partial x^2}+
	\frac{\partial^2 T}{\partial y^2}+
	\frac{\partial^2 T}{\partial z^2}
	\label{eq:159.7}
\end{equation}

в циліндричній системі координат, якщо тепло поширюється тільки в радіальному напрямку:

\begin{equation}
	\nabla^2 T =
	\frac{1}{r}\frac{\partial}{\partial r}
	\left( r\frac{\partial T}{\partial r} \right)
	\label{eq:159.8}
\end{equation}

і в сферичній системі, якщо тепловий потік має сферичну симетрію:

\begin{equation}
	\nabla^2 T =
	\frac{1}{r^2}\frac{\partial}{\partial r}
	\left( r^2\frac{\partial T}{\partial r} \right)
	\label{eq:159.9}
\end{equation}

%---------------------------------------------------------
\begin{figure}[h!]\centering
	\includegraphics[width=1\linewidth]{Pic8}
	\caption{Елементарний об’єм і потоки, що через нього проходять в евклідовій (а), циліндричній (б) та сферичній системах координат (в).}
	\label{pic:Pic8}
\end{figure}
%---------------------------------------------------------

Для прикладу, розглянемо потік тепла, який виходить від дроту, довжина якого набагато більше його діаметру. Кількість тепла, що виходить в одиницю часу через циліндричну поверхню радіусом~$r$ всередині дроту становить $J(r)\,2\pi r L\,dt$ (Рис.~\ref{pic:Pic9}). Через циліндр радіусом $r+dr$ проходить відповідно $J(r+dr)\,2\pi(r+dr)L\,dt$, де $L$ --- довжина дроту. Кількість тепла, що накопичується між циліндрами становить
\[
	Q = \rho c_V\,\frac{\partial T}{\partial t}\,2\pi r L\,dr.
\]

Прирівнюючи отримані величини, отримаємо:

\begin{equation}
	\frac{\partial}{\partial r}\left( J(r)\,2\pi r L \right)
	=
	-\,\rho c_V\,2\pi r L \frac{\partial T}{\partial t}
	\label{eq:159.10}
\end{equation}

Звідки отримуємо (підставляючи $J=-\lambda\,\dfrac{\partial T}{\partial r}$):

\begin{equation}
	\frac{1}{r}
	\frac{d}{dr}
	\left( r\lambda\frac{dT}{dr} \right)
	=
	\rho c_V\,\frac{\partial T}{\partial t}
	\label{eq:159.11}
\end{equation}

В стаціонарному випадку $dT/dt = 0$, отже:

\begin{equation}
	\frac{1}{r}
	\frac{d}{dr}
	\left( r\frac{dT}{dr} \right) = 0
	\label{eq:159.12}
\end{equation}

Інтегруючи:

\begin{equation}
	\frac{dT}{dr} = \frac{A}{r}
	\label{eq:159.13}
\end{equation}

а отже температура:

\begin{equation}
	T(r) = A\ln r + B
	\label{eq:159.14}
\end{equation}

Стала $A$ визначається тепловим потоком від дроту:

\[
	Q = -\lambda\,2\pi L r\,\frac{dT}{dr}
	= -\lambda\,2\pi L A
\]

тому

\[
	A = -\frac{Q}{2\pi L \lambda}.
\]

Звідси для коефіцієнта теплопровідності остаточно:

\begin{equation}
	\lambda =
	\frac{Q}{2\pi L (T_1 - T_2)}
	\ln\frac{r_2}{r_1}
	\label{eq:159.15}
\end{equation}

%---------------------------------------------------------
\begin{figure}[h!]\centering
	\includegraphics[width=0.5\linewidth]{Pic9}
	\caption{До виводу рівняння теплового балансу для нагрітого дроту.}
	\label{pic:Pic9}
\end{figure}
%---------------------------------------------------------

%% --------------------------------------------------------
\section{Умова \texorpdfstring{$\lambda \gg l$}{lambda >> l}. Еффузія}
%% --------------------------------------------------------

Розберемося тепер з явищами в газах, коли довжина вільного пробігу набагато
більше ніж характерні розміри системи і виконується умова $\lambda \gg L$
або хоча б $\lambda \sim L$. Кажуть, що якщо довжина вільного пробігу
молекул в газі одного порядку з розміром посудини, в якому він міститься,
то газ знаходиться в стані вакууму. При цьому якщо $\lambda \gg L$,
то вакуум називають глибоким, $\lambda \ge L$ --- середнім,
і $\lambda \le L$ --- низьким або форвакуумом.

Основна відмінність цього випадку від вже розглянутого раніше, полягає
в тому, що енергія, імпульс та інші фізичні величини передаються молекулами
від однієї частини посудини до іншої безпосередньо, а не через серію
проміжних зіткнень з іншими молекулами. В якомусь сенсі поняття локального
термодинамічної рівноваги втрачає свою актуальність.

У цьому випадку знову стає доречним чисто механістичний підхід, оскільки
зіткненнями молекул газу можна знехтувати. Для зменшення цієї різниці
розглянемо два гази при різних температурах і тисках (в загальному випадку)
по різні боки перегородки. Якщо перегородку прибрати, то умова відсутності
конвекційних потоків буде $p_1 = p_2$ і $T_1 = T_2$ --- умова
термодинамічної рівноваги.

У разі ж, якщо в перегородці зробити малий отвір розміром $l$ такий, щоб
для довжини вільного пробігу молекул $\lambda \gg l$, то в стаціонарних
умовах необхідно, щоб потік частинок зліва був рівним потоку частинок
справа. В силу того, що $\lambda \gg l$ --- розміру отвору, будь-яка частинка,
що пролітає через нього характеризується усередненими параметрами руху,
обумовленими макроскопічними параметрами газу в тій області звідки вона
прилітає, але ніяк не обумовленими зіткненнями молекул в області отвору.

Такий потік молекул через отвір з розмірами багато менше довжини вільного
пробігу називають ефузивним потоком, а саме явище --- ефузією. На практиці
маленький отвір реалізується за допомогою пористої перегородки з розмірами
пор набагато менше ніж $\lambda$.


%% --------------------------------------------------------
\section{Ефект Кнудсена та інші явища ефузії}
%% --------------------------------------------------------


Отже, умова рівноваги між газами $1$ і $2$ в посудинах $А$ і $В$ в загальному випадку записується так (Рис.~\ref{pic:161.1}):
\begin{equation}
	\nu_{1A} S = \nu_{2B} S
	\label{eq:161.1}
\end{equation}
%
де $\nu_i$ --- кількість молекул, які проходять через отвір в одиницю часу.
Скориставшись \eqref{eq:131.3} маємо
\begin{equation}
	\nu_i = \frac{1}{4} n_{iAB} \langle v_{iAB} \rangle
	\label{eq:161.2}
\end{equation}
%
де $n_{iAB} = p_{iAB}/kT_{iAB}$ --- концентрація молекул $i$-того сорту в посудинах $А$ і $В$ відповідно,
$\langle v_{iAB} \rangle$ --- середня швидкість відповідних молекул, що дорівнює $\sqrt{\frac{8kT_{iAB}}{\pi m_i}}$.
Звідки:
\begin{equation}
	n_{1A}\sqrt{\frac{T_{1A}}{m_1}} = n_{2B}\sqrt{\frac{T_{2B}}{m_2}}
	\label{eq:161.3}
\end{equation}
\begin{equation}
	\frac{p_{1A}}{\sqrt{m_1 T_{1A}}}
	=
	\frac{p_{2B}}{\sqrt{m_2 T_{2B}}}
	\label{eq:161.4}
\end{equation}
%
де $S$ --- площа отвору.
Розглянемо декілька окремих випадків.
%---------------------------------------------------------
\begin{figure}[h!]\centering
	\includegraphics[width=0.5\linewidth]{Pic11}
	\caption{До ілюстрації явищ ефузії (а) та теплового ковзання (б)}
	\label{pic:161.1}
\end{figure}
%---------------------------------------------------------
Нехай з різного боку перегородки знаходяться однакові гази $m_i = m$.
Тоді з \eqref{eq:161.3} і \eqref{eq:161.4} досить розглянути одне.
В умовах рівноваги
\begin{equation}
	\frac{p_{1A}}{\sqrt{T_{1A}}}
	=
	\frac{p_{2B}}{\sqrt{T_{2B}}}
	\label{eq:161.5}
\end{equation}
Тобто, якщо спочатку тиски були рівні, то газ починає перетікати з посудини з більш низькою температурою до посудини з більш високою температурою, поки не встановиться співвідношення \eqref{eq:161.5}.
Явище отримало назву \emph{теплової ефузії} або \emph{ефекту Кнудсена}.
Нехай тепер спочатку тиск різних газів і їх температури по різні боки перегородки рівні.
Тоді, в початковий момент часу $p_{1B} = p_{2A}$.
Звідки:
\begin{equation}
	\frac{\nu_{1A}}{\nu_{2B}} = \sqrt{\frac{m_2}{m_1}}
	\label{eq:161.6}
\end{equation}
%
тобто ефузивні потоки через пористу перегородку обернено пропорційні кореню із молекулярної маси газу.
На цьому явищі, яке отримало назву \emph{ізотермічної ефузії}, заснований один з методів розділення ізотопів.
Особливості теплопередачі в розріджених газах обумовлюються також і \emph{радіометричні ефекти}.
Нехай деяка пласка поверхня нагріта нерівномірно так, що температура її зростає в напрямку $x$ (Рис.~\ref{pic:161.1}б).
Молекули газу, які приходять зліва, мають в середньому меншу швидкість, ніж ті, що приходять справа, оскільки температура газу зліва менше, ніж справа.
Це означає, що на тіло діє тангенціальна сила, спрямована справа наліво.
Відповідно на газ діє протилежна сила $F$, яка спричиняє рух газу в напрямку зростання температури.
Це явище отримало назву \emph{теплового ковзання}.
З явищем теплового ковзання пов'язаний і \emph{радіометричний ефек}т.
Він полягає в тому, що нерівномірно нагріті тіла, поміщені в розріджені гази, самовільно починають рухатися в напрямку від більш нагрітого боку до менш нагрітого.
Нерівномірне нагрівання зазвичай пов'язане з нерівномірним освітленням, що і обумовлює назву ефекту.
Внесок у цей ефект приносять дві складові.
Перша --- обумовлена тепловим переміщенням газу від холодних ділянок тіла до гарячих.
Завдяки в’язкості в рух втягується все більша частина газу поблизу тіла.
Внаслідок виконання закону збереження імпульсу тіло починає рухатися в протилежному напрямку.
Швидкість руху $v$ складає:
\begin{equation}
	v = \frac{\lambda}{p} \frac{d\langle v\rangle}{dx}
	\label{eq:161.7}
\end{equation}
%
де $\lambda$ --- довжина вільного пробігу, а $d\langle v\rangle/dx$ --- швидкість зміни середньої теплової швидкості.
Оскільки $\lambda \sim 1/p$, то ефект стає помітний у дуже розріджених газах, хоча не надто сильно розріджених.
Зокрема, він проявляється в осіданні пилу за радіаторами центрального опалення.
Друга складова радіометричного ефекту обумовлена тим, що нормальна складова швидкості молекул, що відлітають від більш нагрітої ділянки, більша, ніж від менш нагрітої.
Ця складова є основною у сильно розріджених газах і пропорційна тиску газу.


%% --------------------------------------------------------
\section{Абсолютний манометр Кнудсена}
%% --------------------------------------------------------

На радіометричному ефекті ґрунтується дія \emph{манометра Кнудсена} або \emph{радіометричного вакуумметра}, що складається з двох пластин нагрітих до різних температур.

Прилад містить нерухому пластину (температура \(T_1\)) і рухому (або підвішену систему), нагріту до \(T_2 > T_1\). Молекули, що відбиваються від гарячої поверхні, передають більший імпульс, ніж від холодної.

Точний розрахунок Кнудсена (з урахуванням повної теплової акомодації) дає силу
\begin{equation}
	F = \frac{1}{2} p S \left( \sqrt{\frac{T_2}{T_1}} - 1 \right)
\end{equation}
і момент крутіння (в оригінальній конструкції)
\begin{equation}
	M = p S l \left( \sqrt{\frac{T_2}{T_1}} - 1 \right),
\end{equation}
де \(l\) --- плече сили, \(S\) ---- площа пластин.

\begin{figure}[h!]
	\centering
	\includegraphics[width=0.4\linewidth]{Pic12}
	\caption{Схема радіометричного манометра Кнудсена}
	\label{pic:Pic12}
\end{figure}

До переваг такого манометра можна віднести те, що він виконує абсолютні вимірювання, які не залежність від складу газу (при \(\alpha \approx 1\)).
Недоліками є чутливість до \(\Delta T\), обмежений діапазон тисків та теплова інерція.

Сучасні комерційні прилади (Knudsen gauge, радіометричні вакуумметри) зберігають той самий фізичний принцип, хоча геометрія може бути іншою (мембрана, кварцевий резонатор тощо).

\Closesolutionfile{answer}